\chapter{課題リポジトリを取得する}

課題ファイルは GitHub 上で管理されています。
\cmd{git clone} コマンドで WSL のホームディレクトリに取得します。
\textbf{この操作は最初の1回だけ}行います。

% =================================================
\section{git clone で取得する}
% =================================================

WSL ターミナルで以下を実行します。

\begin{wslbox}[リポジトリのクローン]
\begin{lstlisting}[style=bash]
$ cd ~
$ git clone \
  https://github.com/shiru2/NMR-lab-kadai.git
\end{lstlisting}
\end{wslbox}

正常に完了すると \cmd{NMR-lab-kadai} フォルダが作成されます。

\begin{wslbox}[取得の確認]
\begin{lstlisting}[style=bash]
$ ls \
  ~/NMR-lab-kadai/2026年度_B4課題引継ぎ/
No.1_FID  No.2_FFT1D  No.3_FFT2D
No.4_DFT_exercise  No.5_CompressedSensing
No.6_Denoising  No.7_Metrics  No.8_CT
README.md  pyproject.toml  uv.lock
\end{lstlisting}
\end{wslbox}

\cmd{pyproject.toml} と \cmd{uv.lock} が表示されれば正しく取得できています。

\begin{notebox}[git が使えない場合]
大学のネットワーク環境によって \cmd{git clone} が
ブロックされる場合があります。
その場合は GitHub のページから zip をダウンロードして展開してください。\\[0.5em]
\texttt{Code} $\rightarrow$ \texttt{Download ZIP} を選び、
展開したフォルダを\\
\cmd{{\textbackslash}{\textbackslash}wsl.localhost{\textbackslash}Ubuntu{\textbackslash}home{\textbackslash}username}\\
に配置してください（\cmd{username} は自分の WSL ユーザー名）。
\end{notebox}

% =================================================
\section{2回目以降：最新版に更新する}
% =================================================

課題ファイルが更新された場合は \cmd{git pull} で最新版を取得できます。

\begin{wslbox}[最新版に更新]
\begin{lstlisting}[style=bash]
$ cd ~/NMR-lab-kadai
$ git pull
\end{lstlisting}
\end{wslbox}
